% =====================================================================
%  CV — Huawei Canada, Noah's Ark Laboratory, Markham
%        Co-op Researcher : trois requisitions couvertes par ce seul CV —
%        « Video Generation & 3D Computer Vision », « Machine Learning and
%        Computer Vision », « Computer Vision and Machine Learning ».
%  Dérivé de master-cv.md (source de vérité unique)
%
%  CIBLAGE — différences avec la version Graphics Technology Lab :
%   - Les trois annonces demandent « video / image understanding and
%     synthesis » → la synthèse remonte (VQ-VAE, render-then-estimate),
%     le vocabulaire purement graphique (SVBRDF, neural rendering) descend.
%   - « Excellent coding skills with Python and/or C++ » → C et C++ sont
%     listés SANS gras. Le relevé porte un 7/20 à l'examen de C++ (S7) :
%     mettre C++ en avant sur un CV accompagné du relevé se retourne contre.
%   - « OpenCV […] is an asset » → OpenCV devient visible, adossé à la
%     bibliothèque d'augmentation écrite de zéro chez Cali.
%   - « Familiarity with Linux » et « linear algebra, probability theory,
%     optimization » → ligne dédiée, ce sont leurs mots.
%   - « Academic research experience, for example participation in paper
%     publishing, is an asset » → les publications restent en premier.
%
%  CE QUI DESCEND : rendu de maillages, masques de silhouette, calibration.
%  Noah's Ark n'est pas un labo graphique.
%
%  CE QUI N'EST PAS REVENDIQUÉ, faute de preuve : Stable Diffusion, Flux,
%  GAN, modèles de diffusion, VLM, développement mobile. Les annonces les
%  citent comme atouts ; les inventer se paierait au premier entretien.
%
%  INTERDITS (master-cv.md §11) : projet clical, −58 %, PFA 82,3 %,
%  audience Medium, Google Final Round, projets perso, noms de clients,
%  rejet antérieur du papier ICLR, canal Talents Handicap, mandarin
%  « working knowledge » (niveau réel : débutant).
%
%  ⚠ AVANT ENVOI — deux points à trancher, voir la conversation :
%    (1) ces co-ops sont annoncés 8 à 16 mois ; la disponibilité réelle est
%        février → octobre 2027, soit 8 mois. Le formulaire ne doit cocher
%        que « 8 months », jamais 12 ni 16.
%    (2) autorisation de travail au Canada pour 2027 : rien n'est écrit ici
%        tant que le statut n'est pas établi.
% =====================================================================

\documentclass[10pt, letterpaper]{article}

\usepackage[
    ignoreheadfoot,
    top=0.7 cm,
    bottom=0.55 cm,
    left=1.15 cm,
    right=1.15 cm,
    footskip=0.8 cm,
]{geometry}
\usepackage{titlesec}
\usepackage{enumitem}
\usepackage{fontawesome5}
\usepackage[dvipsnames]{xcolor}
\usepackage{textcomp}

\definecolor{MyRed}{RGB}{160, 20, 40}

\usepackage[
    pdftitle={Melissa Colin — CV — Huawei Canada, Co-op Researcher, Noah's Ark Lab},
    pdfauthor={Melissa Colin},
    colorlinks=true,
    urlcolor=MyRed
]{hyperref}

\usepackage{charter}

\pagestyle{empty}
\setlength{\parindent}{0pt}

\titleformat{\section}
  {\normalsize\bfseries\scshape\raggedright\color{MyRed}}{}{0em}{}[\titlerule]
\titlespacing{\section}{0pt}{6pt}{4pt}

\setlist[itemize]{leftmargin=*, noitemsep, topsep=1pt, partopsep=0pt, parsep=0pt}

\newcommand{\entry}[4]{%
  \noindent \textbf{#1} \hfill #2 \\
  \textit{#3} \hfill \textit{#4} \par
}

\IfFileExists{personal.tex}{\input{personal.tex}}{}
\providecommand{\myphone}{+33 7 82 52 XX XX}

\begin{document}

% ---------------------------------------------------------------- HEADER
\begin{center}
    {\fontsize{19}{19}\selectfont \textbf{Mélissa COLIN}} \\[3pt]
    {\normalsize Computer Vision \& Deep Learning $\cdot$ 3D Human Motion, Video Understanding and Synthesis} \\[3pt]

    \small{
    \faMapMarker* Bordeaux, France \;|\;
    \faPhone* \myphone{} \;|\;
    \faEnvelope~\href{mailto:contact@melissacolin.ai}{contact@melissacolin.ai} \;|\;
    \faGlobe~\href{https://melissacolin.ai}{melissacolin.ai} \;|\;
    \faGithub~\href{https://github.com/melissa-colin}{melissa-colin} \;|\;
    \faGraduationCap~\href{https://scholar.google.com/citations?user=7r7iFpsAAAAJ}{Scholar}
    }
\end{center}

\vspace{1pt}

{\small \textit{Final-year MEng student in computer science, first author on a CVPR 2027
submission. I design neural architectures for 3D human motion from video, and the pipelines
that synthesise and estimate their training data at corpus scale. Expected graduation 09/27,
available from February 2027 for a co-op of up to eight months, through October 2027.}}

% ---------------------------------------------------------------- PUBLICATIONS
\section{Publications}
\begin{itemize}
    \item \textbf{Colin, M.}, Wang, Y., Cheng, L. \textit{InterForce: Physically Admissible Multi-Person Motion Refinement.} First author, in preparation, targeting \textbf{CVPR 2027}.
    \item \textbf{Colin, M.}, Chraibi Kaadoud, I. \textit{Performances and Explainability of ViT and CNN Architectures: An Empirical Study Using LIME, SHAP and Grad-CAM.} RJCIA / PFIA 2024, La Rochelle. First author, published, presented orally. \href{https://hal.science/hal-04641791}{[HAL]}
    \item \textbf{Co-author}, 3D group dance generation. In preparation, targeting ICLR 2027.
\end{itemize}

% ---------------------------------------------------------------- RESEARCH
\section{Research Experience}

\entry{Vision and Learning Lab, University of Alberta (Prof.\ Li Cheng)}{May 2026 -- Sep 2026}%
{Visiting Research Student, \textbf{Mitacs Globalink Research Award}}{Edmonton, Canada}
\begin{itemize}
    \item Designed \textbf{two neural architectures}: a permutation-equivariant transformer whose attention is factorised over people, joints and time, so a scene is handled without ever fixing the number of people, and a second model whose output is forces, not poses.
    \item Wrote a \textbf{differentiable rigid-body dynamics engine from scratch}: recursive Newton--Euler inverse dynamics, composite-rigid-body mass matrix, implicit actuation, joint limits. Checked against nine analytic invariants to machine precision, and differentiated through inside the training loop.
    \item Built a \textbf{27{,}000-clip paired corpus by render-then-estimate}: clean motion is synthesised into video, then re-estimated from those pixels, and the estimation error becomes the paired noisy sample. Multi-person video, camera calibration, occlusion ordering.
    \item \textbf{Dense monocular depth} (\textit{Depth Anything}) brought to metric scale by anchoring it on two quantities measurable directly in the image, then used to correct depth ordering between people: \textbf{$-$33\% inter-person penetration over 100\% of the corpus}, purely algorithmic, nothing trained, no ground truth.
\end{itemize}

\vspace{1pt}

\entry{ENSEIRB-MATMECA --- Research Initiation, 17/20}{Feb -- May 2026}%
{Reproduction and critical study of \textit{InterMask} (ICLR 2025)}{Bordeaux, France}
\begin{itemize}
    \item Reproduced the \textbf{VQ-VAE} baseline (FID 1.034 $\to$ \textbf{0.918}), then tested two hypotheses on the physical constraints of the loss. Neither improved FID or R-Precision; \textbf{reported the negative result} instead of dropping it.
\end{itemize}

% ---------------------------------------------------------------- INDUSTRY
\section{Industry Experience}

\entry{Sector Group}{Jun -- Aug 2025}{R\&D Engineer, AI and Safety --- information extraction from \textbf{PDFs with no text layer}, which ruled out parsing and forced vision and document processing together}{Paris, France}

\vspace{1pt}

\entry{Cali Intelligences (deeptech start-up, real-time video analytics)}{Jan 2023 -- Sep 2024}%
{Applied AI Intern, then AI Developer (apprenticeship)}{Talence, France}
\begin{itemize}
    \item Took detection accuracy from \textbf{60\% to 90\%} by designing and building the entire training pipeline (YOLOv7 $\to$ YOLOv8): ingestion, label conversion, an image augmentation library written from scratch in OpenCV, pre-training validation, and YAML-driven orchestration.
    \item Shipped it to production on Docker, GitHub Actions and K3s.
\end{itemize}

% ---------------------------------------------------------------- EDUCATION
\section{Education}

\entry{ENSEIRB-MATMECA, Bordeaux INP}{Sep 2024 -- Expected 09/27}%
{MEng in Computer Science (\textit{Diplôme d'ingénieur}), AI track, final year}{Bordeaux, France}
\begin{itemize}
    \item Ranked \textbf{5\textsuperscript{th}/81 (top 6\%)} in semester 8 and \textbf{3\textsuperscript{rd}/93} in semester 6 (French /20 scale, transcript available). \textbf{Deep Learning 18/20} $\cdot$ \textbf{Research Initiation 17/20}.
\end{itemize}

\vspace{1pt}

\entry{EPSI Bordeaux}{Sep 2021 -- Sep 2024}%
{BSc in Artificial Intelligence \& Data Science, ranked \textbf{1\textsuperscript{st}/10} in final year}{Bordeaux, France}

% ---------------------------------------------------------------- SKILLS
\section{Technical Skills}
\begin{itemize}
    \item \textbf{Deep learning:} \textbf{PyTorch} (advanced) $\cdot$ Keras/TensorFlow $\cdot$ Transformers and ViT $\cdot$ \textbf{VQ-VAE} $\cdot$ CNNs $\cdot$ LoRA/PEFT $\cdot$ explainability: LIME, SHAP, Grad-CAM $\cdot$ written from scratch: backpropagation, PCA, K-Means, logistic regression.
    \item \textbf{Computer vision \& 3D:} \textbf{OpenCV} $\cdot$ video and image synthesis $\cdot$ monocular depth $\cdot$ SMPL / SMPL-X, forward kinematics, axis-angle and 6D rotations $\cdot$ camera calibration, occlusion ordering $\cdot$ differentiable rigid-body dynamics.
    \item \textbf{Programming \& foundations:} \textbf{Python} $\cdot$ C $\cdot$ C++ $\cdot$ Java $\cdot$ JavaScript $\cdot$ SQL $\cdot$ Bash $\cdot$ Linux/POSIX, Git, Docker, multi-GPU clusters $\cdot$ numerical linear algebra, probability, optimisation, NP-completeness.
    \item \textbf{Languages:} French (native) $\cdot$ English \textbf{B2, TOEIC 840/990}, placement conducted entirely in English $\cdot$ Mandarin (elementary).
\end{itemize}

\end{document}
